% !TEX root = ../main.tex
%-----------------------------------------------------------------------
%  The naive scaling table, and the two caveats that stop it being read
%  as a cost estimate.
%  \input by ../main.tex -- not compiled on its own.
%-----------------------------------------------------------------------
\hd{Naive scaling.}
Arithmetic, not measurement. A $\times4$ refinement in each horizontal direction
with the vertical left alone --- itself a guess about how the grid would actually
be built, and one of the first things a real design would revisit.

\smallskip
\begin{tabularx}{\textwidth}{@{}lY{1.0}Y{1.0}@{}}
\toprule
 & \textbf{1\dg\ (\code{DINO\_1deg})} & \textbf{1/4\dg, if scaled $\times4$} \\
\midrule
Grid
  & $51 \times 198 \times 36 = 363{,}528$
  & $204 \times 792 \times 36 \approx 5.82$\,M cells ($16\times$) \\
\addlinespace
Spacing
  & 1\dg\ zonal; meridional follows $\cos(\mathrm{lat})$,
    ${\sim}0.35\dg$ to $1.0\dg$, mean ${\sim}0.71\dg$
  & a quarter of each \\
\addlinespace
Timestep
  & \code{dT = 1800 s}
  & ${\sim}450$\,s if the advective CFL sets it ($4\times$ the steps) \\
\addlinespace
Work per model year
  & ---
  & $\mathbf{\approx 64\times}$ \\
\addlinespace
Decomposition
  & \code{sNx=17, sNy=22, nPx=3, nPy=9} $\rightarrow$ 27 ranks
  & to be chosen; \code{-n} and \code{SIZE.h} move together \\
\bottomrule
\end{tabularx}

\hd{What the $\times64$ applies to.}
The measured 1\dg\ figures are a 200-year spin-up at 32\,h\,21\,m for
3{,}513{,}600 forward steps on 27 ranks, and a five-year reference adjoint at
23\,h\,03\,m. The $\times64$ is \emph{work}, not wall time --- the point of a
bigger grid is that it also takes many more ranks, so the wall-clock figure
depends entirely on a decomposition nobody has chosen yet.

\hd{Do not scale the 23-hour figure without reading it first.}
\code{useGrdchk = .TRUE.} in both setups, so that run also paid for the built-in
gradient check's perturbed forward integrations. With the check off the same
adjoint is about 9 hours. Scaling 23 hours rather than 9 overstates the 1/4\dg\
adjoint by roughly a factor of 2.5, on the strength of a one-line namelist
setting --- which is the kind of error that decides a direction is unaffordable
before anyone has checked what was being measured.

\hd{The spin-up is the item to look at first.}
A 200-year spin-up is what makes the 1\dg\ work possible, and at 1/4\dg\ it is
the dominant cost of the whole direction by a wide margin --- far more than the
adjoint the direction actually exists to run. Whether it can be shortened,
initialised from the 1\dg\ state, or avoided altogether is the first thing worth
deciding, and it should be decided before any of the questions in section~1 are
costed.
